\par\noindent\textbf{(1)} 因 \([\alpha a^\dagger,-\alpha^*a]=|\alpha|^2\) 为常数，BCH公式给 \(D(\alpha)=e^{-|\alpha|^2/2}e^{\alpha a^\dagger}e^{-\alpha^*a}\)。施于真空得 \(|\alpha\rangle=e^{-|\alpha|^2/2}\sum_{n=0}^\infty\alpha^n|n\rangle/\sqrt{n!}\)。
\par\noindent\textbf{(2)} 光子数概率 \(P_n=|\langle n|\alpha\rangle|^2=e^{-|\alpha|^2}|\alpha|^{2n}/n!\)，为Poisson分布。故 \(\bar n=|\alpha|^2\)、\(\operatorname{Var}(n)=|\alpha|^2\)，Fano因子 \(F=\operatorname{Var}(n)/\bar n=1\)，不同于热态的超Poisson涨落。
\par\noindent\textbf{(3)} 用 \(x=\sqrt{\hbar/(2m\omega)}(a+a^\dagger)\)、\(p=-i\sqrt{\hbar m\omega/2}(a-a^\dagger)\) 及 \(a|\alpha\rangle=\alpha|\alpha\rangle\)，得 \((\Delta x)^2=\hbar/(2m\omega)\)、\((\Delta p)^2=\hbar m\omega/2\)，从而 \(\Delta x\Delta p=\hbar/2\) 恰好饱和。
\par\noindent\textbf{(4)} 在 \(H=\hbar\omega(a^\dagger a+1/2)\) 下，\(e^{-iHt/\hbar}|\alpha\rangle=e^{-i\omega t/2}|\alpha e^{-i\omega t}\rangle\)。因此 \(\langle x(t)\rangle\) 和 \(\langle p(t)\rangle\) 沿经典谐振子椭圆轨道旋转，而波包方差保持基态宽度。
